%% ACM SIGGRAPH paper (conference track) on Derive.
%%
%% The class is acmart, which Overleaf and TeX Live ship (SIGGRAPH needs acmart 2.16 or
%% newer; CTAN carries 2.20). For the journal track (ACM Transactions on Graphics) change
%% the first line to \documentclass[acmtog]{acmart}; add `anonymous,review` for a
%% double-blind submission.
%%
%% derive.sty is Derive's package for tables and figures that update without a new
%% version: \derivetable{name} and \derivefigure{name}. The rendered page shows each
%% slot's current value, and "Download LaTeX source" writes the current data next to this
%% file as derive-dynamic/<name>.tex, so the paper compiles with what the page shows.
\documentclass[sigconf]{acmart}
\usepackage{derive}

%% SIGGRAPH requires author-year citations.
\citestyle{acmauthoryear}

%% Rights and venue. TAPS fills the DOI and ISBN at production; leave the placeholders.
\setcopyright{acmlicensed}
\copyrightyear{2026}
\acmYear{2026}
\acmDOI{XXXXXXX.XXXXXXX}
\acmConference[SIGGRAPH '26]{Special Interest Group on Computer Graphics and Interactive Techniques Conference}{July 2026}{Vancouver, BC, Canada}
\acmISBN{978-1-4503-XXXX-X/26/07}

\begin{document}

\title{Paper Title}

\author{First Author}
\email{first@example.edu}
\orcid{0000-0000-0000-0000}
\affiliation{%
  \institution{Institution}
  \city{City}
  \country{Country}
}

\author{Second Author}
\email{second@example.org}
\affiliation{%
  \institution{Institution}
  \city{City}
  \country{Country}
}

\begin{abstract}
One paragraph: the problem, why it matters, what this paper does, and what the reader
gains. Replace this text.
\end{abstract}

%% The CCS concepts block comes from https://dl.acm.org/ccs; paste the generated XML
%% here and mirror it with \ccsdesc lines below.
\begin{CCSXML}
<ccs2012>
   <concept>
       <concept_id>10010147.10010371.10010372</concept_id>
       <concept_desc>Computing methodologies~Rendering</concept_desc>
       <concept_significance>500</concept_significance>
   </concept>
</ccs2012>
\end{CCSXML}

\ccsdesc[500]{Computing methodologies~Rendering}

\keywords{keyword one, keyword two, keyword three}

%% The teaser is a dynamic figure: replace its image from the Data tab (or PATCH the
%% `teaser` slot) and the page and the export follow, without a new version.
\begin{teaserfigure}
  \derivefigure[width=\textwidth]{teaser}
  \caption{Teaser caption: what the reader sees at a glance.}
  \Description{Describe the teaser for readers who cannot see it.}
  \label{fig:teaser}
\end{teaserfigure}

\maketitle

\section{Introduction}
\label{sec:intro}

State the problem and why existing work leaves it open~\cite{example2020method}. Say
what this paper contributes, in the order the sections deliver it. Figure~\ref{fig:teaser}
should already make the point.

\section{Related Work}

Group prior work by the approach it takes rather than by year. Cite with
\verb|\cite{key}| for a parenthetical reference~\cite{example2023survey} and
\verb|\citet{key}| when the authors are the subject of the sentence, as
\citet{example2020method} did.

\section{Method}
\label{sec:method}

Introduce notation once. A displayed equation gets a number when it is referenced:
\begin{equation}
  L(\mathbf{x}, \omega_o) = L_e(\mathbf{x}, \omega_o) + \int_{\Omega} f_r(\mathbf{x}, \omega_i, \omega_o)\, L_i(\mathbf{x}, \omega_i)\, (\omega_i \cdot \mathbf{n})\, d\omega_i .
  \label{eq:rendering}
\end{equation}
Equation~\eqref{eq:rendering} is the rendering equation; inline math such as
$\omega_i \cdot \mathbf{n}$ reads well in running text.

\section{Results}
\label{sec:results}

Table~\ref{tab:results} is a dynamic table: an agent (or you) fills it in through
Derive's data API as runs finish, and every reader sees the current numbers.

\begin{table}
  \caption{Quantitative comparison. Higher is better.}
  \label{tab:results}
  \derivetable{results}
\end{table}

A static table works too, with booktabs rules:

\begin{table}
  \caption{A static table.}
  \label{tab:static}
  \begin{tabular}{lrr}
    \toprule
    Method & PSNR & SSIM \\
    \midrule
    Baseline & 27.1 & 0.91 \\
    Ours & \textbf{28.6} & \textbf{0.94} \\
    \bottomrule
  \end{tabular}
\end{table}

\section{Conclusion}

Close with what was shown, what it does not cover, and what comes next.

\begin{acks}
Thank the people and grants that made the work possible. This section is dropped under
the \texttt{anonymous} option.
\end{acks}

\bibliographystyle{ACM-Reference-Format}
\bibliography{references}

\end{document}
